G\aa rding's theory establishes the convexity of hyperbolicity
cones~\cite{garding-1959}. The original Lax conjecture asks for a
definite determinantal representation of homogeneous hyperbolic polynomials in three
variables~\cite{lax-1958}. Helton and Vinnikov proved the corresponding
definite determinantal representation theorem~\cite{helton-vinnikov-2007};
Lewis, Parrilo, and Ramana established its equivalence with Lax's
formulation~\cite{lewis-parrilo-ramana-2005}. In higher dimension, Br\"and\'en
constructed a real-zero polynomial none of whose positive powers admits
a definite determinantal representation~\cite[Theorem~3.3]{branden-2011}.
That obstruction does not exclude a different determinant, with additional
factors, from defining the same cone. The geometric conjecture concerns this remaining
possibility. Kummer exhibited the distinction concretely: the cone of
the specialized V\'amos polynomial has a $7\times7$ representation,
although no positive power of that polynomial has a definite
determinantal representation~\cite[Section~3]{kummer-2016-vamos}.
Raghavendra, Ryder, Srivastava, and Weitz proved exponential
lower bounds on the matrix size of spectrahedral representations for
suitable hyperbolicity cones~\cite{raghavendra-ryder-srivastava-2018}.
Recent positive results include Netzer's theorem for
hyperbolic cubics in five variables~\cite[Theorem~3.9]{netzer-2026}.

Kummer and Netzer prove spectrahedrality for cones of strictly
hyperbolic polynomials~\cite[Theorem~1]{kummer-netzer-2026-strict}.
Here \emph{strictly hyperbolic} means that $p(te-x)$ has distinct roots
whenever $x\notin\R e$. The polynomial~\eqref{eq:polynomial} fails this
hypothesis: at $w=(I_4,0,0)\notin\R e$ it gives
\[
 p(te-w)=(t-1)^{16}t^4.
\]
For the reduced polynomial of Corollary~\ref{cor:reduced}, the same line
gives $q(te-w)=(t-1)^{12}t^4$. Thus neither defining polynomial satisfies
the strict-case hypothesis.

A \emph{spectrahedral shadow} is a linear image of a set defined by an
affine real symmetric linear matrix inequality.
Netzer and Sanyal proved that a hyperbolicity cone is a spectrahedral
shadow when every nonzero boundary point is a smooth point of the
defining polynomial~\cite[Theorem~1.1]{netzer-sanyal-2015}.
Our result concerns a direct representation in the original variables;
it makes no assertion about representations using auxiliary variables.

The scalar ingredient is the coefficient-one Choi--Lam biquadratic
form~\cite[Section~4]{choi-lam-1977}, which reduces the coefficient of
the extra cyclic squared monomials from two in Choi's earlier
nonnegative form~\cite{choi-1975} to one. Its inability to dominate a nonzero
bilinear square follows from the extremality proved in
\cite[Theorem~4.4]{choi-lam-1977}; we give a short proof of precisely this
property. It is called \emph{weak extremality} by Blekherman, Rai\c{t}\u{a},
Shankar, and Sinn~\cite[Definition~2.7]{blekherman-raita-shankar-sinn-2022},
whose characterization uses Taylor expansions of dominated squares at
zeros~\cite[Theorem~3.2]{blekherman-raita-shankar-sinn-2022}.
The correspondence between nonnegative biquadratic forms and
positive maps on real symmetric matrices is classical; see
\cite[Section~3]{ha-2013}. Positive maps have also been used to obtain
certificates of conic polynomial stability~\cite{dey-gardoll-theobald-2021}.
Here the construction turns the scalar obstruction into an obstruction
to every finite pencil describing one hyperbolicity cone. A related
general strategy appears in Scheiderer's local sums-of-squares
obstructions to spectrahedral shadows~\cite[Section~4]{scheiderer-2018}.
Our proof extracts bilinear square minorants from a differentiated
operator-norm identity; it does not invoke a criterion for shadows.

Gonz\'alez Nevado states a positive solution of the geometric conjecture
in~\cite[Theorem~53, pp.~16--17]{gonzalez-nevado-2026}.
Appendix~\ref{app:opposition} records a counterexample to the
path-containment assertion used in that argument. The construction
and proof in the present paper are independent of that comparison.
